% =====================================================================
% Next -- minimal revision snippets (EN for main/supplement, CN for reading edition)
% Packages assumed: tabularx, array.  No journal-class-specific commands.
% Define once in the preamble (used by the two wide tables):
%   \newcolumntype{Y}{>{\raggedright\arraybackslash}X}
% Every number below is taken from the supplied packet (README, analysis.json,
% audit.json, rotating-quorum-regression.log, version-bridge.{md,json}).
% =====================================================================
 
% ---------------------------------------------------------------------
% [1] MAIN TEXT, directly after Table II (Sec. III-E).  EN: 82 words.
% ---------------------------------------------------------------------
% --- EN ---
\paragraph{Terminal holders and Intent.}
A verified TXCer alone creates no redeemable public obligation. If its source
does not execute and the holder cannot obtain a valid successor that executes
publicly, the protocol offers no direct redemption or unconditional exit, and
waiting does not create one. Intent is a global deduplication choice: for one
Network, Subject and Nonce, the ledger executes at most one TxID.
Organizations do not coordinate it during certification, so conflicting
payments can both obtain QCs; a QC fixes content and liability, not execution.
 
% --- CN ---
\paragraph{终端持证人与 Intent。}
仅持有已验证的 TXCer 不产生可请求兑付的公开义务。若其来源交易不执行，且持证人无法获得并公开执行一个合法后继，
协议不提供直接兑付或无条件退出路径，等待本身也不会产生该权利。Intent 是全局去重语义的选择：对同一 Network、
Subject 与 Nonce，账本至多执行一个 TxID。各组织在认证阶段不预先协调该选择，因此相互冲突的支付可以各自获得有效 QC；
QC 固定内容与责任，而不保证执行。
 
% ---------------------------------------------------------------------
% [2] SUPPLEMENT: member approval vs public execution (new table).
% ---------------------------------------------------------------------
% --- EN ---
\begin{table*}[!t]
\caption{Member approval versus public execution. D: deterministic functions of the
submitted bytes and static configuration. S: checks against two different states.
L: delivery and scheduling assumptions. Theorem~3 compares only histories in which every
payment and increase succeeds; the table does not claim that a QC-bearing payment
always executes.}
\label{tab:sup-checks}
\footnotesize
\renewcommand{\arraystretch}{1.15}
\begin{tabularx}{\textwidth}{@{}>{\raggedright\arraybackslash}p{2.6cm}YYY@{}}
\hline
Check & Member approval (\texttt{ApproveDirectBytes}) & Public execution (\texttt{VerifyDirectSubmission}, \texttt{EvaluateDirectPaymentAt}) & Relation and claim \\
\hline
\multicolumn{4}{@{}l}{\textit{D. Deterministic}} \\
D1 Encoding, owner authorization, initial funding openings, rule IDs, network, configuration, epoch
 & \texttt{PrepareDirectVector}; identity equals own organization
 & Same \texttt{PrepareDirectVector}; organization looked up by configuration
 & Same function on the same bytes and static configuration. \\
D2 Admission vector and caps (CAL cap includes change; escrow floor includes one compensation charge per certificate input)
 & \texttt{PrepareDirectVector}
 & \texttt{PrepareDirectVector}
 & Identical caps on both sides. \\
D3 Direct input certificates; value conservation
 & Issuer configuration, rules and output checked; one certificate per certificate input, no duplicate or unused certificate; $\sum$in $=\sum$out
 & The same checks, plus network equality; $\sum$in $=\sum$out at evaluation
 & Equivalent checks; the member also lies in a validated single-network configuration. \\
\hline
\multicolumn{4}{@{}l}{\textit{S. State}} \\
S1 Input instances
 & Per-instance candidate/consumed lock and local creation record (derived from authenticated blocks), written before signing
 & Consumed flag; final creation record, or a direct certificate for a missing output
 & Lemma~2: no two conflicting QCs. A QC excludes public consumption by a different fact. Member lag can refuse, not mis-approve. \\
S2 Intent
 & Organization-local Intent $\to$ TxID
 & Global Intent $\to$ TxID, written on successful execution
 & Independent states; conflict is possible with valid QCs (Lemma~5). \\
S3 Grants, capacity, fees
 & Worker-slice \texttt{Available} $\ge$ cap for every resource; fee inputs and grants checked in member state
 & CAL: registration within the grant (\texttt{Reserved}+\texttt{Spent} $\le B_g$); non-CAL: \texttt{updateUsage}; FUEL account or fee-input UTXOs
 & CAL: Theorem~1 and Eq.~(7) exclude a CAL-capacity or reserve failure at first registration. Non-CAL resources, FUEL balance and fee-input spendability: no corresponding claim. \\
\hline
\multicolumn{4}{@{}l}{\textit{L. Liveness}} \\
L1 Inclusion, delivery, resubmission
 & Original signer keeps the approval and follows the chain
 & Honest inclusion and fair delivery; no global Intent conflict for the source
 & No completion bound and no timeout unlocking; adaptation is not required for closure. \\
\hline
\end{tabularx}
\end{table*}
 
% --- CN ---
\begin{table*}[!t]
\caption{成员批准与公开执行的对照。D：只取决于提交字节与静态配置的确定性函数；S：对两个不同状态的检查；
L：传送与调度假设。定理~3 仅比较每笔支付与增资都成功的历史；本表不主张带 QC 的支付一定执行。}
\label{tab:sup-checks-cn}
\footnotesize
\renewcommand{\arraystretch}{1.15}
\begin{tabularx}{\textwidth}{@{}>{\raggedright\arraybackslash}p{2.6cm}YYY@{}}
\hline
检查项 & 成员批准（\texttt{ApproveDirectBytes}） & 公开执行（\texttt{VerifyDirectSubmission}、\texttt{EvaluateDirectPaymentAt}） & 关系与主张 \\
\hline
\multicolumn{4}{@{}l}{\textit{D. 确定性检查}} \\
D1 编码、所有者授权、初始资金开口、规则标识、网络/配置/纪元
 & \texttt{PrepareDirectVector}；身份等于本组织
 & 同一 \texttt{PrepareDirectVector}；按配置查组织
 & 相同字节与静态配置下是同一函数。 \\
D2 准入向量与额度（CAL 额度含找零；托管下限含每个证书输入一次赔付费用）
 & \texttt{PrepareDirectVector}
 & \texttt{PrepareDirectVector}
 & 两侧额度相同。 \\
D3 直接输入证书；价值守恒
 & 核对发行组织配置、规则与输出；每个证书输入恰一个证书，无重复、无多余证书；$\sum$输入 $=\sum$输出
 & 同样检查，另加网络相等；评估时 $\sum$输入 $=\sum$输出
 & 等价检查；成员处于已验证的单网络配置内。 \\
\hline
\multicolumn{4}{@{}l}{\textit{S. 状态检查}} \\
S1 输入实例
 & 逐实例候选/已消费锁与本地创建记录（来自已认证区块），签名前写入
 & 已消费标记；最终创建记录，或对缺失输出的直接证书
 & 引理~2：不存在两个冲突 QC；QC 排除不同事实对其输入的公开消费。成员落后只会拒绝，不会误批。 \\
S2 Intent
 & 组织本地 Intent $\to$ TxID
 & 全局 Intent $\to$ TxID，成功执行时写入
 & 状态相互独立；可在两方都有效 QC 时冲突（引理~5）。 \\
S3 授权、额度、费用
 & 各资源的 Worker 切片 \texttt{Available} $\ge$ 额度；费用输入与授权按成员状态检查
 & CAL：登记不超授权（\texttt{Reserved}+\texttt{Spent} $\le B_g$）；非 CAL：\texttt{updateUsage}；FUEL 账户或费用输入 UTXO
 & CAL：定理~1 与式~(7) 排除首次登记时 CAL 额度或储备不足。非 CAL 资源、FUEL 余额与费用输入可花费性：无对应主张。 \\
\hline
\multicolumn{4}{@{}l}{\textit{L. 活性}} \\
L1 纳入、传送、补交
 & 原签署者保留批准并跟随区块链
 & 诚实纳入与公平传送；来源无全局 Intent 冲突
 & 无完成时限，不以超时解锁；闭合不需要适配。 \\
\hline
\end{tabularx}
\end{table*}
 
% ---------------------------------------------------------------------
% [3] SUPPLEMENT S4: automatic source resubmission (R).
% ---------------------------------------------------------------------
% --- EN ---
\paragraph{Automatic source resubmission (R).}
On the frozen build, one run schedules 2,000 ordinary payments (100/s for 20\,s)
plus three cross-organization parent--child pairs, 2,006 payments in total, with
automatic resubmission enabled and adaptation not paused. For each pair, the driver
obtains the parent's three signatures through a collector that neither persists,
installs nor submits it, and the driver never submits the parent itself; the collector
stays alive and observes, so this is withholding, not process-crash injection. The
child spends the parent's TXCer. After the child's public execution exposes that
certificate, the parent's original signers rebuild and submit the parent. All three
obligations end Fulfilled and none is compensated; all 2,006 payments close,
committee state hashes agree, outboxes are empty, and member CAL, execution and
storage reservations return to zero. The three pairs are three recovery targets,
not repetitions. Times are client observations; outbox creation times were not
recorded.
 
\begin{table}[!t]
\caption{Pair-level observations in R (ms). Receipt times start at each payment's own
first send; the last two columns start at the child's first send and are client
observations, not internal recovery times.}
\label{tab:sup-r}
\centering\footnotesize
\begin{tabular}{@{}rrrrr@{}}
\hline
Pair & Parent & Child & Child & Source \\
 & receipt & receipt & public & Fulfilled \\
\hline
0 & 100.146 & 116.559 & 1391.849 & 2767.211 \\
1 &  82.988 &  61.749 & 1806.036 & 2958.003 \\
2 &  75.033 & 135.119 & 1286.147 & 2442.890 \\
\hline
\end{tabular}
\end{table}
 
% --- CN ---
\paragraph{来源自动补交（R）。}
在冻结构建上，一轮计划 2,000 笔普通付款（100/s，20\,s）与三个跨组织父子对，共 2,006 笔，
自动补交开启、适配不暂停。每个父子对中，驱动通过一个既不持久化、不 INSTALL、也不提交的收集器取得父交易三个签名，
驱动自身从不提交父交易；收集器保持在线并继续观察，因此这是扣留发送，不是进程崩溃注入。子交易花费父交易的 TXCer；
子交易公开执行并暴露该证书后，父交易的原签署成员重建并提交父交易。三个义务均以 Fulfilled 结束，无一赔付；
2,006 笔全部闭合，委员会状态哈希一致，outbox 清空，成员 CAL、执行与存储预留归零。三个父子对是三个恢复目标，
不是重复实验。时间为客户端观察值；未记录 outbox 的内部创建时间。
 
\begin{table}[!t]
\caption{R 中各父子对的观察值（ms）。到账时间自各笔付款自身首次发送起算；
后两列自子交易首次发送起算，为客户端观察点，不是内部恢复时间。}
\label{tab:sup-r-cn}
\centering\footnotesize
\begin{tabular}{@{}rrrrr@{}}
\hline
父子对 & 父到账 & 子到账 & 子公开 & 来源 Fulfilled \\
\hline
0 & 100.146 & 116.559 & 1391.849 & 2767.211 \\
1 &  82.988 &  61.749 & 1806.036 & 2958.003 \\
2 &  75.033 & 135.119 & 1286.147 & 2442.890 \\
\hline
\end{tabular}
\end{table}
 
% ---------------------------------------------------------------------
% [3b] SUPPLEMENT S3: rotating-signer boundary (one paragraph).
% ---------------------------------------------------------------------
% --- EN ---
\paragraph{Rotating-signer bound.}
A deterministic regression uses production approvals, public execution and
compensation. One Byzantine member signs without applying its local limit; the three
honest members co-sign in the pairs $(0,1)$, $(0,2)$, $(1,2)$. Each round certifies a
100-CAL source that loses on global Intent; a valid child then consumes its output and
the issuer is compensated. With $B_g=300$ the reserve falls $300\to200\to100\to0$,
honest reservations are $(100,100,0)$, $(200,100,100)$, $(200,200,200)$, and the fourth
request is refused by all three honest members. A repeated decision changes nothing and
CAL/FUEL are conserved. The run attains $Q_g=B_g$, so the two-round stop in
Table~S4 is specific to a fixed signer set. It measures neither attacker profit nor
cost: the source inputs stay locked and FUEL is spent.
 
% --- CN ---
\paragraph{轮换签署者边界。}
一个确定性回归使用生产批准、公开执行与赔付入口。一个拜占庭成员绕过自身本地额度签名；三个诚实成员按
$(0,1)$、$(0,2)$、$(1,2)$ 轮换共同签署。每轮认证一个 100 CAL、在全局 Intent 上落败的来源，
再由合法子交易消费其输出并赔付发行方。取 $B_g=300$，备付依次 $300\to200\to100\to0$，诚实成员累计预留为
$(100,100,0)$、$(200,100,100)$、$(200,200,200)$，第四笔请求被三个诚实成员全部拒绝。重复决定无任何变化，
CAL/FUEL 守恒。该运行达到 $Q_g=B_g$，故表~S4 中两轮即止只对固定签署者集合成立。
它既不度量攻击者收益，也不度量成本：来源输入仍被锁定，FUEL 已花费。
 
% ---------------------------------------------------------------------
% [4] CONCLUSION (compressed).  EN: see word count in reply.
% ---------------------------------------------------------------------
% --- EN ---
Next separates continuation, economic closure and historical record. On one host
with synchronous member commits, certified continuation reaches the last receipt of a
100-hop chain in 4.707\,s, against 64.479\,s when each hop waits for public execution
(retries included). Reserve-backed signing capacity covers certificate liability,
original signers can resubmit exposed sources, and a due compensation decision closes
an obligation that a late source repays, even with adaptation paused. Optional
representation records the reserve debit under the original block identity. Terminal
holders have no direct redemption; open deployment accepts bounded reserve loss.
 
% --- CN ---
Next 将续付、经济闭合与历史记录分开。在单机、成员同步提交的配置下，凭证续付使 100 跳链的最后一次收据在
4.707\,s 到达，逐跳等待公开执行为 64.479\,s（均含重试）。以备付支撑的签名额度覆盖证书责任，原签署成员可补交已暴露的来源，
到期的赔付决定可闭合义务并由迟到来源归还，适配暂停时亦然。可选的授权表示把备付借记记录在原区块身份之下。
终端持证人没有直接兑付路径；开放部署接受有界的备付损失。
 
% ---------------------------------------------------------------------
% [5] SUPPLEMENT: version bridge and wait-retry wording.
% ---------------------------------------------------------------------
% --- EN ---
\paragraph{Version bridge.}
Between 721800c and 9879f64, five member production files changed
(\texttt{blocks.go}, \texttt{direct.go}, \texttt{outcome.go}, \texttt{partial\_reclaim.go},
\texttt{source\_recovery.go}); the 721800c N/R/C/P and capital results therefore keep their
original build and are not inherited. Between 9879f64 and f856c7b, the checked member,
committee, rules, redaction, transport, wallet, protocol, finality and Comet production
files are unchanged; only client and test files differ. The three redaction core files and
the listed Comet overlay files have identical SHA-256 digests in all three builds. No
older number is relabelled.
 
\paragraph{Retries in wait mode.}
Wait mode uses 414 attempts for 300 payments (140, 138 and 136 per run). The client
sleeps 50\,ms after every failed attempt, whatever the cause (transport, HTTP status or
read error), and the original logs do not classify failures. The 114 extra attempts
therefore include at least 5.7\,s of fixed sleep over the three runs, plus unmeasured
service time. We report the $13.70\times$ ratio with retries included and derive no
retry-free figure.
 
% --- CN ---
\paragraph{版本桥接。}
721800c 至 9879f64 之间，五个成员生产文件有变动（\texttt{blocks.go}、\texttt{direct.go}、\texttt{outcome.go}、
\texttt{partial\_reclaim.go}、\texttt{source\_recovery.go}）；因此 721800c 的 N/R/C/P 与资本结果保留其原构建，不被继承。
9879f64 至 f856c7b 之间，所检查的成员、委员会、规则、表示、传输、钱包、协议、终局与 Comet 生产文件无变化，
仅客户端与测试文件不同。三个表示核心文件及所列 Comet overlay 文件在三个构建中的 SHA-256 完全相同。旧数字均不改标。
 
\paragraph{等待模式的重试。}
等待模式 300 笔付款共 414 次尝试（每轮 140、138、136）。客户端在每次失败尝试后固定休眠 50\,ms，
不论原因是传输、HTTP 状态还是读取错误；原日志没有逐次分类。因此 114 次额外尝试在三轮中至少含 5.7\,s 的固定休眠，
另有未测量的服务时间。我们报告含重试的 $13.70\times$，不推导无重试数字。
 